\chapter{Gebräuchliche Grenzwerte und Mittelwerte}
\label{b21guide:ch:standard}

In diesem Kapitel beginnen Folgen mit Wurzeln oder Nenner $n$
bei $n=1$. Beliebige Ergänzungen bei $n=0$ ändern ihre Grenzwerte
nicht. Potenzen mit natürlichem Exponenten und reelle Wurzeln
werden im Sinn der Arithmetik reeller Zahlen verwendet.

\section{Vergleichsfolgen}

\begin{theorem}[Kehrwerte und feste Potenzen natürlicher Zahlen]
Für jede feste ganze Zahl $p\geq1$ gilt
\[
 n^p\to+\infty,\qquad \frac1{n^p}\to0.
\]
\end{theorem}

\begin{theorem}[Geometrische Folgen]
\label{b21guide:thm:geometric}
Für $q\in\mathbb R$ hat die Folge $(q^n)_{n\geq1}$
folgendes Verhalten:
\begin{enumerate}
 \item $|q|<1$: Es gilt $q^n\to0$.
 \item $q=1$: Die Folge ist konstant gleich $1$.
 \item $q=-1$: Die Folge ist divergent und hat genau die
 Häufungswerte $-1$ und $1$.
 \item $q>1$: Es gilt $q^n\to+\infty$.
 \item $q<-1$: Es gilt $|q^n|\to+\infty$; die geraden
 Potenzen streben gegen $+\infty$, die ungeraden gegen $-\infty$.
\end{enumerate}
\end{theorem}

\begin{theorem}[Wurzeln mit wachsendem Index]
\label{b21guide:thm:nth-roots}
Für jedes feste $c>0$ gilt
\[
 \sqrt[n]{c}\to1,\qquad \sqrt[n]{n}\to1.
\]
\end{theorem}

\begin{theorem}[Geometrischer Abfall gegenüber Potenzwachstum]
Für jede feste ganze Zahl $p\geq0$ und $0<q<1$ gilt
\[
 n^p q^n\to0.
\]
Entsprechend gilt für $c>1$ die Konvergenz $n^p/c^n\to0$.
\end{theorem}

\begin{theorem}[Die Eulersche Folge]
\label{b21guide:thm:euler}
Die Folge
\[
 e_n=\left(1+\frac1n\right)^n\qquad(n\geq1)
\]
ist streng wachsend und nach oben beschränkt.
Ihr reeller Grenzwert wird mit $e$ bezeichnet.
Es gilt $2<e<3$. Außerdem ist
\[
 f_n=\left(1+\frac1n\right)^{n+1}
\]
streng fallend, $e_n<e<f_n$ für alle $n\geq1$ und $f_n\to e$.
\end{theorem}

\begin{theorem}[Fakultätswachstum im Vergleich]
\label{b21guide:thm:factorial-growth}
Für $n\geq1$ sei $n!=1\cdot2\cdots n$.
Für jede feste reelle Zahl $c$ gilt
\[
 \frac{c^n}{n!}\longrightarrow0.
\]
Außerdem gelten
\[
 \frac{n!}{n^n}\longrightarrow0,\qquad
 \frac{\sqrt[n]{n!}}{n}\longrightarrow\frac1e.
\]
Hier ist $e$ der Grenzwert der Eulerschen Folge aus
Satz~\ref{b21guide:thm:euler}.
\end{theorem}

\section{Subadditives und multiplikatives Wachstum}

\begin{definition}[Sub- und Superadditivität; multiplikative Fassungen]
\label{b21guide:def:subadditive}
Eine bei $1$ beginnende reelle Folge $(a_n)$ heißt
\emph{subadditiv}, wenn
\[
 a_{m+n}\leq a_m+a_n\qquad(m,n\geq1),
\]
und \emph{superadditiv}, wenn stets die umgekehrte Ungleichung gilt.

Eine positive reelle Folge $(u_n)_{n\geq1}$ heißt
\emph{submultiplikativ}, wenn
\[
 u_{m+n}\leq u_mu_n\qquad(m,n\geq1),
\]
und \emph{supermultiplikativ}, wenn stets die umgekehrte
Ungleichung gilt. Alle vier Bedingungen betreffen die
Verknüpfung der \emph{Indizes}; Monotonie wird nicht verlangt.
\end{definition}

\begin{theorem}[Feketes Grenzwertsatz]
\label{b21guide:thm:fekete}
Für Folgen mit den Eigenschaften aus
Definition~\ref{b21guide:def:subadditive} gelten:
\begin{enumerate}
 \item Ist $(a_n)$ subadditiv, so konvergiert $(a_n/n)$
 gegen eine reelle Zahl oder strebt gegen $-\infty$.
 Der Grenzwert ist genau dann reell, wenn es ein $c\in\mathbb R$
 mit $a_n\geq cn$ für alle $n\geq1$ gibt.
 \item Ist $(a_n)$ superadditiv, so konvergiert $(a_n/n)$
 gegen eine reelle Zahl oder strebt gegen $+\infty$.
 Der Grenzwert ist genau dann reell, wenn es ein $C\in\mathbb R$
 mit $a_n\leq Cn$ für alle $n\geq1$ gibt.
 \item Ist $(u_n)$ positiv und submultiplikativ, so
 konvergiert $(\sqrt[n]{u_n})$ gegen eine reelle Zahl
 $\lambda\geq0$. Der Grenzwert $0$ ist zugelassen.
 \item Ist $(u_n)$ positiv und supermultiplikativ, so
 konvergiert $(\sqrt[n]{u_n})$ gegen eine reelle Zahl
 $\lambda>0$ oder strebt gegen $+\infty$.
\end{enumerate}
Weder die Quotientenfolge $(a_n/n)$ noch die Wurzelfolge
$(\sqrt[n]{u_n})$ muß monoton sein.
Ein Grenzwert von $u_{n+1}/u_n$ wird nicht vorausgesetzt.
\end{theorem}

\section{Arithmetische Mittel}

\begin{definition}[Cesàro-Mittel einer Folge]
\label{b21guide:def:cesaro}
Für eine reelle Folge $(a_n)_{n\in\mathbb N}$ heißt
\[
 \sigma_n=\frac1{n+1}\sum_{k=0}^{n}a_k
\]
die Folge ihrer \emph{arithmetischen Mittel} oder
\emph{Cesàro-Mittel}. Die Summe ist stets endlich.
\end{definition}

\begin{theorem}[Cesàroscher Mittelwertsatz]
\label{b21guide:thm:cesaro}
Aus $a_n\to L\in\mathbb R$ folgt $\sigma_n\to L$.
Die Umkehrung gilt im allgemeinen nicht.
Auch bestimmte Divergenz wird erhalten:
\[
 a_n\to+\infty\ \Longrightarrow\ \sigma_n\to+\infty,
 \qquad
 a_n\to-\infty\ \Longrightarrow\ \sigma_n\to-\infty.
\]
\end{theorem}

\begin{theorem}[Taubersches Kriterium für Cesàro-Mittel]
\label{b21guide:thm:cesaro-converse}
Für eine reelle Folge $(a_n)$ gebe es $C\geq0$ und $N\geq1$ mit
\[
 n|a_n-a_{n-1}|\leq C\qquad(n\geq N).
\]
Dann folgt aus der Konvergenz ihrer Cesàro-Mittel
$\sigma_n=(n+1)^{-1}\sum_{k=0}^{n}a_k$ bereits
\[
 \sigma_n\to L\in\mathbb R\quad\Longrightarrow\quad a_n\to L.
\]
Auch $\sigma_n\to+\infty$ beziehungsweise $\sigma_n\to-\infty$
erzwingt unter dieser Voraussetzung die entsprechende bestimmte
Divergenz von $(a_n)$.

Gilt die stärkere Bedingung $n(a_n-a_{n-1})\to0$, so gilt sogar
$a_n-\sigma_n\to0$, ohne daß zuvor ein Grenzwert der Mittel
vorausgesetzt werden muß.
Die schwächere Bedingung $a_n-a_{n-1}\to0$ allein ersetzt
die angegebene Voraussetzung nicht.
\end{theorem}

\begin{theorem}[Stolz--Cesàro-Kriterium]
\label{b21guide:thm:stolz}
Seien $(a_n),(b_n)$ reelle Folgen, wobei $(b_n)$ strikt wachsend
ist und $b_n\to+\infty$ gilt.
Für $L\in\mathbb R$ gilt auf einer Restfolge
\[
 \frac{a_{n+1}-a_n}{b_{n+1}-b_n}\longrightarrow L
 \quad\Longrightarrow\quad
 \frac{a_n}{b_n}\longrightarrow L.
\]
Die Differenznenner sind nach Voraussetzung positiv,
und $b_n$ ist schließlich positiv.
Monotonie von $(a_n)$ wird nicht verlangt.
Strebt der Differenzenquotient gegen $+\infty$ beziehungsweise
$-\infty$, so strebt auch $a_n/b_n$ in dieselbe Richtung.
\end{theorem}

\begin{theorem}[Stolz--Cesàro in der Nullfolgenfassung]
\label{b21guide:thm:stolz-zero}
Seien $a_n\to0$ und $b_n\to0$, wobei $b_n>0$ und
$b_{n+1}<b_n$ für alle $n$ gelten.
Für $L\in\mathbb R$ gilt
\[
 \frac{a_{n+1}-a_n}{b_{n+1}-b_n}\to L
 \quad\Longrightarrow\quad
 \frac{a_n}{b_n}\to L.
\]
Der Zähler darf sein Vorzeichen wechseln.
Die Bedingung $a_n\to0$ ist in dieser Fassung erforderlich.
Unter denselben Voraussetzungen gilt die Aussage auch für
bestimmte Divergenz des Differenzenquotienten gegen
$+\infty$ oder $-\infty$; dann hat $a_n/b_n$ dieselbe
Divergenzrichtung.
\end{theorem}

\begin{theorem}[Asymptotische Steigung]
Ist $a_{n+1}-a_n\to L\in\mathbb R$, so gilt
\[
 \frac{a_n}{n}\longrightarrow L\qquad(n\geq1).
\]
Insbesondere folgt aus verschwindenden benachbarten Differenzen
$a_n/n\to0$; daraus folgt keine Konvergenz von $(a_n)$ selbst.
\end{theorem}

\section{Wurzelgrenzwerte und weitere Mittelwerte}

\begin{theorem}[Wurzelkriterium für Folgen]
\label{b21guide:thm:root-test}
Für eine reelle Folge $(a_n)_{n\geq1}$ existiere
\[
 \sqrt[n]{|a_n|}\longrightarrow q\in[0,+\infty).
\]
Ist $q<1$, so gilt $a_n\to0$.
Ist $q>1$, so gilt $|a_n|\to+\infty$.
Für $q=1$ folgt keine Entscheidung über Konvergenz oder Divergenz.
Anders als beim Quotientenkriterium sind hier auch Nullglieder erlaubt.
\end{theorem}

\begin{theorem}[Vom Quotienten- zum Wurzelgrenzwert]
\label{b21guide:thm:ratio-root}
Für eine positive reelle Folge $(a_n)_{n\geq1}$ gilt
\[
 \frac{a_{n+1}}{a_n}\longrightarrow q\in[0,+\infty)
 \quad\Longrightarrow\quad
 \sqrt[n]{a_n}\longrightarrow q.
\]
Auch aus $a_{n+1}/a_n\to+\infty$ folgt
$\sqrt[n]{a_n}\to+\infty$.
Die Umkehrung wird durch die Existenz des Wurzelgrenzwerts
allein nicht gewährleistet.
\end{theorem}

\begin{definition}[Gewichtete arithmetische Mittel]
\label{b21guide:def:weighted-means}
Für reelle $a_k$ und Gewichte $w_k\geq0$ setzen wir
\[
 W_n=\sum_{k=0}^{n}w_k,\qquad
 \mu_n=\frac{\sum_{k=0}^{n}w_ka_k}{W_n}
 \quad\text{an den Indizes mit }W_n>0.
\]
\end{definition}

\begin{theorem}[Gewichteter Cesàro-Satz]
\label{b21guide:thm:weighted-cesaro}
Es seien $w_k\geq0$ und $W_n\to+\infty$.
Dann sind die gewichteten Mittel $\mu_n$ schließlich definiert.
Es gilt für jedes $L\in\mathbb R$
\[
 a_n\to L\quad\Longrightarrow\quad\mu_n\to L.
\]
Auch aus $a_n\to+\infty$ beziehungsweise $a_n\to-\infty$
folgt $\mu_n\to+\infty$ beziehungsweise $\mu_n\to-\infty$.
Einzelne Gewichte dürfen verschwinden.
\end{theorem}

\begin{definition}[Geometrische und harmonische Mittel]
\label{b21guide:def:geometric-harmonic}
Für eine positive reelle Folge $(a_n)_{n\in\mathbb N}$ seien
\[
 G_n=\sqrt[n+1]{\prod_{k=0}^{n}a_k},
 \qquad
 J_n=\frac{n+1}{\sum_{k=0}^{n}1/a_k}.
\]
Die Produkte und Summen sind endlich. $G_n$ heißt das
\emph{geometrische}, $J_n$ das \emph{harmonische Mittel}.
\end{definition}

\begin{theorem}[Grenzwerte geometrischer und harmonischer Mittel]
\label{b21guide:thm:geometric-harmonic}
Sei $a_n>0$ für alle $n$.
Aus $a_n\to L$ mit $L\geq0$ folgt
\[
 G_n\to L,\qquad J_n\to L.
\]
Der Grenzwert $L=0$ ist ausdrücklich eingeschlossen.
Aus $a_n\to+\infty$ folgt auch $G_n\to+\infty$ und $J_n\to+\infty$.
\end{theorem}

\section{Grenzwerterhaltende endliche Mittelungsschemata}

\begin{definition}[Endliches lineares Mittelungsschema]
\label{b21guide:def:regular-scheme}
Für jedes $n\in\mathbb N$ seien reelle Zahlen
$w_{n,0},\ldots,w_{n,n}$ vorgegeben.
Zu einer reellen Folge $(a_n)$ gehört dann die Folge
\[
 t_n=\sum_{k=0}^{n}w_{n,k}a_k.
\]
Das Schema heißt \emph{regulär}, wenn für jede reelle Folge
$a_n\to L\in\mathbb R$ auch $t_n\to L$ gilt.
Die Gewichte dürfen negativ sein und mit der Zeile $n$ wechseln.
Alle auftretenden Summen sind endlich.
\end{definition}

\begin{theorem}[Silverman--Toeplitz für endliche Zeilen]
\label{b21guide:thm:toeplitz}
Ein endliches lineares Mittelungsschema ist genau dann regulär,
wenn die folgenden drei Bedingungen gleichzeitig erfüllt sind:
\begin{enumerate}
 \item Es gibt $C\geq0$ mit
 $\sum_{k=0}^{n}|w_{n,k}|\leq C$ für alle $n$.
 \item Für jedes feste $k\in\mathbb N$ gilt $w_{n,k}\to0$
 beim Grenzübergang $n\to\infty$ mit $n\geq k$.
 \item Es gilt $\sum_{k=0}^{n}w_{n,k}\to1$.
\end{enumerate}
Sind alle Gewichte nichtnegativ und summieren sie sich in
jeder Zeile zu $1$, so ist das Schema genau dann regulär,
wenn für jedes feste $k$ die Gewichte $w_{n,k}$ gegen $0$ gehen.
\end{theorem}

